% GENERATED FILE. Edit canonical lessons, then run npm run generate:books.
% canonical-source-hash: fnv1a64:cdbd415c1fc8ff2c
% canonical-lessons: MR-C309-gamavne, MR-C309-sosne, MR-C309-mrtyu, MR-C309-maran-pavne, MR-C309-udbhavne

\chapter{Losses}
\label{ch:mr-losses}
\hlchaptermodality{309}

\begin{chapteropening}
\textbf{By the end of this chapter:} \emph{I can talk about loss and death in a news story.}

Complete the last lesson of chapter 309: I can talk about loss and death in a news story.
\end{chapteropening}

\section[\texorpdfstring{\mr{गमावणे} (gamāvṇe)}{गमावणे (gamāvṇe)}]{\mr{गमावणे} (gamāvṇe) — to lose forever}

\label{lesson:MR-C309-gamavne}



\begin{quote}
Before the new one: say the Marathi for a lineage, then the Marathi for to preserve.
\end{quote}

\subsection*{You'll want to know: \mr{गमावणे}}
\textbf{\mr{गमावणे}} — \emph{gamāvṇe} — \textquotedblleft{}to lose forever\textquotedblright{}.

\begin{quote}
\mr{त्या} \mr{पुरात} \mr{अनेकांनी} \mr{आपली} \mr{घरं} \mr{गमावली}.
\end{quote}

\begin{quote}
\emph{tyā purāt anekānnī āplī gharaṁ gamāvlī.}
\end{quote}

\begin{quote}
Many people lost their homes in that flood.
\end{quote}

Past, with -\mr{ने} (he …): \mr{त्याने} \mr{गमावलं}, \emph{tyāne gamāvlaṁ}.

For what does not come back: a life, a chance, a home. \mr{हरवणे} is for mislaying.

\begin{cousinweb}
\textbf{\mr{गमावणे}} (\emph{gamāvṇe}) is a tadbhava: CDIAL lists Marathi \emph{gamāvṇe}, 'to lose', under Sanskrit \emph{gamyate}, 'one goes' (CDIAL 4033): to let go away.
\end{cousinweb}

\begin{grammarlens}[title={What you've built}]
One more word for telling what happened.
\end{grammarlens}

\subsection*{Your turn}
\begin{itemize}
\raggedright
  \item \emph{Say it:} \emph{gamāvṇe}
  \item \emph{Say it:} \emph{gamāvṇe}, once more
  \item \emph{Say it:} \emph{tyā purāt anekānnī āplī gharaṁ gamāvlī.} — Many people lost their homes in that flood.
  \item \emph{Recall:} say the Marathi for a lineage, then the Marathi for to preserve, then say \emph{gamāvṇe} again
\end{itemize}

\begin{tcolorbox}[breakable,colback=teal!4,colframe=teal!35!black,title={Before you move on}]
What does \mr{गमावणे} mean? (\textquotedblleft{}To lose forever\textquotedblright{}.) Say it once more.
\end{tcolorbox}
\section[\texorpdfstring{\mr{सोसणे} (sosṇe)}{सोसणे (sosṇe)}]{\mr{सोसणे} (sosṇe) — to endure; to bear}

\label{lesson:MR-C309-sosne}



\begin{quote}
Before the new one: say the Marathi for to preserve, then the Marathi for to lose forever.
\end{quote}

\subsection*{You'll want to know: \mr{सोसणे}}
\textbf{\mr{सोसणे}} — \emph{sosṇe} — \textquotedblleft{}to endure; to bear\textquotedblright{}.

\begin{quote}
\mr{त्या} \mr{कुटुंबाने} \mr{खूप} \mr{दुःख} \mr{सोसलं}.
\end{quote}

\begin{quote}
\emph{tyā kuṭumbāne khūp duḥkh soslaṁ.}
\end{quote}

\begin{quote}
That family endured a lot of sorrow.
\end{quote}

Past, with -\mr{ने} (he …): \mr{त्याने} \mr{सोसलं}, \emph{tyāne soslaṁ}.

Transitive: \mr{दुःख} \mr{सोसणे}, to bear grief; \mr{सहन} \mr{करणे} is the formal cousin.

\begin{cousinweb}
\textbf{\mr{सोसणे}} (\emph{sosṇe}) is a Marathi verb of bearing pain; its origin is not settled.
\end{cousinweb}

\begin{grammarlens}[title={What you've built}]
One more word for telling what happened.
\end{grammarlens}

\subsection*{Your turn}
\begin{itemize}
\raggedright
  \item \emph{Say it:} \emph{sosṇe}
  \item \emph{Say it:} \emph{sosṇe}, once more
  \item \emph{Say it:} \emph{tyā kuṭumbāne khūp duḥkh soslaṁ.} — That family endured a lot of sorrow.
  \item \emph{Recall:} say the Marathi for to preserve, then the Marathi for to lose forever, then say \emph{sosṇe} again
\end{itemize}

\begin{tcolorbox}[breakable,colback=teal!4,colframe=teal!35!black,title={Before you move on}]
What does \mr{सोसणे} mean? (\textquotedblleft{}To endure; to bear\textquotedblright{}.) Say it once more.
\end{tcolorbox}
\section[\texorpdfstring{\mr{मृत्यू} (mṛtyū)}{मृत्यू (mṛtyū)}]{\mr{मृत्यू} (mṛtyū) — death}

\label{lesson:MR-C309-mrtyu}



\begin{quote}
Before the new one: say the Marathi for to lose forever, then the Marathi for to endure.
\end{quote}

\subsection*{You'll want to know: \mr{मृत्यू}}
\textbf{\mr{मृत्यू}} — \emph{mṛtyū} — \textquotedblleft{}death\textquotedblright{}.

\begin{quote}
\mr{त्या} \mr{अपघातात} \mr{दोघांचा} \mr{मृत्यू} \mr{झाला}.
\end{quote}

\begin{quote}
\emph{tyā apghātāt doghāñcā mṛtyū jhālā.}
\end{quote}

\begin{quote}
Two people died in that accident.
\end{quote}

Masculine.

X-\mr{चा} \mr{मृत्यू} \mr{होणे}, X dies.

\begin{cousinweb}
\textbf{\mr{मृत्यू}} (\emph{mṛtyū}) is a tatsama, taken from Sanskrit \emph{mṛtyu}, 'death' (Monier-Williams).
\end{cousinweb}

\begin{grammarlens}[title={What you've built}]
One more word for telling what happened.
\end{grammarlens}

\subsection*{Your turn}
\begin{itemize}
\raggedright
  \item \emph{Say it:} \emph{mṛtyū}
  \item \emph{Say it:} \emph{mṛtyū}, once more
  \item \emph{Say it:} \emph{tyā apghātāt doghāñcā mṛtyū jhālā.} — Two people died in that accident.
  \item \emph{Recall:} say the Marathi for to lose forever, then the Marathi for to endure, then say \emph{mṛtyū} again
\end{itemize}

\begin{tcolorbox}[breakable,colback=teal!4,colframe=teal!35!black,title={Before you move on}]
What does \mr{मृत्यू} mean? (\textquotedblleft{}Death\textquotedblright{}.) Say it once more.
\end{tcolorbox}
\section[\texorpdfstring{\mr{मरण} \mr{पावणे} (maraṇ pāvṇe)}{मरण पावणे (maraṇ pāvṇe)}]{\mr{मरण} \mr{पावणे} (maraṇ pāvṇe) — to pass away}

\label{lesson:MR-C309-maran-pavne}



\begin{quote}
Before the new one: say the Marathi for to endure, then the Marathi for death.
\end{quote}

\subsection*{You'll want to know: \mr{मरण} \mr{पावणे}}
\textbf{\mr{मरण} \mr{पावणे}} — \emph{maraṇ pāvṇe} — \textquotedblleft{}to pass away\textquotedblright{}.

\begin{quote}
\mr{ते} \mr{महान} \mr{लेखक} \mr{गतवर्षी} \mr{मरण} \mr{पावले}.
\end{quote}

\begin{quote}
\emph{te mahān lekhak gatvarṣī maraṇ pāvle.}
\end{quote}

\begin{quote}
That great writer passed away last year.
\end{quote}

Past (he …): \mr{तो} \mr{मरण} \mr{पावला}, \emph{to maraṇ pāvlā}.

Formal, for news: X \mr{मरण} \mr{पावला}. \mr{मरणे} is the plain verb.

\begin{cousinweb}
\textbf{\mr{मरण} \mr{पावणे}} (\emph{maraṇ pāvṇe}) is built from \emph{maraṇ}, 'death', from Sanskrit \emph{maraṇa}, 'dying' (CDIAL 9869), and \emph{pāvṇe}, 'to reach, obtain': to reach death.
\end{cousinweb}

\begin{grammarlens}[title={What you've built}]
One more word for telling what happened.
\end{grammarlens}

\subsection*{Your turn}
\begin{itemize}
\raggedright
  \item \emph{Say it:} \emph{maraṇ pāvṇe}
  \item \emph{Say it:} \emph{maraṇ pāvṇe}, once more
  \item \emph{Say it:} \emph{te mahān lekhak gatvarṣī maraṇ pāvle.} — That great writer passed away last year.
  \item \emph{Recall:} say the Marathi for to endure, then the Marathi for death, then say \emph{maraṇ pāvṇe} again
\end{itemize}

\begin{tcolorbox}[breakable,colback=teal!4,colframe=teal!35!black,title={Before you move on}]
What does \mr{मरण} \mr{पावणे} mean? (\textquotedblleft{}To pass away\textquotedblright{}.) Say it once more.
\end{tcolorbox}
\section[\texorpdfstring{\mr{उद्भवणे} (udbhavṇe)}{उद्भवणे (udbhavṇe)}]{\mr{उद्भवणे} (udbhavṇe) — to arise; to crop up}

\label{lesson:MR-C309-udbhavne}



\begin{quote}
Before the new one: say the Marathi for death, then the Marathi for to pass away.
\end{quote}

\subsection*{You'll want to know: \mr{उद्भवणे}}
\textbf{\mr{उद्भवणे}} — \emph{udbhavṇe} — \textquotedblleft{}to arise; to crop up\textquotedblright{}.

\begin{quote}
\mr{पावसामुळे} \mr{पाण्याचा} \mr{नवा} \mr{प्रश्न} \mr{उद्भवला}.
\end{quote}

\begin{quote}
\emph{pāvsāmuḷe pāṇyācā navā praśna udbhavlā.}
\end{quote}

\begin{quote}
A new water problem arose because of the rain.
\end{quote}

Past (he …): \mr{तो} \mr{उद्भवला}, \emph{to udbhavlā}.

Intransitive and formal: \mr{प्रश्न} \mr{उद्भवणे}, a problem arises.

\begin{cousinweb}
\textbf{\mr{उद्भवणे}} (\emph{udbhavṇe}) is built on Sanskrit \emph{udbhava}, 'arising, origin' (Monier-Williams), with the Marathi verb ending \emph{-ṇe}.
\end{cousinweb}

\begin{grammarlens}[title={What you've built}]
One more word for telling what happened.
\end{grammarlens}

\subsection*{Your turn}
\begin{itemize}
\raggedright
  \item \emph{Say it:} \emph{udbhavṇe}
  \item \emph{Say it:} \emph{udbhavṇe}, once more
  \item \emph{Say it:} \emph{pāvsāmuḷe pāṇyācā navā praśna udbhavlā.} — A new water problem arose because of the rain.
  \item \emph{Recall:} say the Marathi for death, then the Marathi for to pass away, then say \emph{udbhavṇe} again
\end{itemize}

\begin{tcolorbox}[breakable,colback=teal!4,colframe=teal!35!black,title={Before you move on}]
What does \mr{उद्भवणे} mean? (\textquotedblleft{}To arise; to crop up\textquotedblright{}.) Say it once more.
\end{tcolorbox}
